\documentclass[t,12pt]{beamer}

% A4 portrait pages.  The global `t` option top-aligns every frame so the
% lower part of each page remains intentionally open instead of overflowing.
\geometry{paperwidth=21cm,paperheight=29.7cm}

% Compile directly in Overleaf with pdfLaTeX.
\usepackage[T1]{fontenc}
\usepackage{lmodern}
\usepackage{amsmath,amssymb,bm}
\usepackage{booktabs}
\usepackage{array}
\usepackage{siunitx}
\usepackage{xcolor}
\DeclareSIUnit{\AUD}{AUD}

\definecolor{bodygreen}{HTML}{405344}
\definecolor{softgreen}{HTML}{EAF0EA}
\definecolor{accentblue}{HTML}{0072BD}
\definecolor{accentorange}{HTML}{D95319}
\definecolor{warningred}{HTML}{B22222}

\setbeamercolor{normal text}{fg=bodygreen,bg=white}
\setbeamercolor{frametitle}{fg=black,bg=white}
\setbeamercolor{structure}{fg=bodygreen}
\setbeamercolor{alerted text}{fg=warningred}
\setbeamerfont{frametitle}{size=\LARGE,series=\bfseries}
\setbeamertemplate{navigation symbols}{}
\setbeamersize{text margin left=1.2cm,text margin right=1.2cm}
\setbeamertemplate{itemize item}{\raise0.15ex\hbox{\scriptsize$\bullet$}}
\setbeamertemplate{itemize subitem}{\raise0.10ex\hbox{\tiny$\circ$}}
\setbeamertemplate{footline}{%
  \hfill\color{gray}\scriptsize Basic 1.3 three-pump MILP\hspace{0.7em}
  \insertframenumber/\inserttotalframenumber\hspace{1.0em}\vspace{0.5em}}

% Normal readable spacing is retained because portrait pages provide enough
% vertical room; content starts at the top and does not fill the lower page.
\setlength{\parskip}{0.35em}
\sisetup{per-mode=symbol,detect-all}

\newcommand{\code}[1]{\texttt{#1}}
\newcommand{\unitcost}{\mathrm{AUD/m^3}}

\title{Basic 1.3 Core Equations}
\subtitle{Three-pump rolling MILP and economic interpretation}
\author{}
\date{}

\begin{document}

% ----------------------------------------------------------------------
\begin{frame}
  \titlepage
  \vspace{-1.4em}
  \begin{center}
    \small
    Fixed-flow synthetic pumps, five-minute control intervals, and a
    24-hour prediction horizon.
  \end{center}
\end{frame}

% ----------------------------------------------------------------------
\begin{frame}{Key updates from Basic 1.2 to Basic 1.3}
\small
\begin{tabular}{>{\raggedright\arraybackslash}p{0.25\textwidth}
                >{\raggedright\arraybackslash}p{0.31\textwidth}
                >{\raggedright\arraybackslash}p{0.34\textwidth}}
\toprule
Component & Basic 1.2 & Basic 1.3 \\
\midrule
Pump model & One fixed-flow pump & Three independent fixed-flow pumps \\
Decision variable & Scalar binary $u(k)$ & Binary vector
  $\bm u(k)=[u_1,u_2,u_3]$ \\
Optimization & Single-pump binary dynamic programming &
  Three-pump MILP solved by \code{intlinprog} \\
Available modes & Off or P1 & All eight three-pump combinations \\
Switching & Change in one pump state & Separate transition variable for
  every pump \\
Controller output & One immediate request & Full $N\times3$ plan plus
  first-step request \\
Execution feedback & One actual state and flow & Three actual states,
  mode, total flow, total power, and tracking error \\
Closed-loop memory & Previous scalar action & Previous
  \emph{actually executed} $1\times3$ pump state \\
\bottomrule
\end{tabular}

\vspace{0.4em}
\small
The outer rolling structure is unchanged: five-minute sampling, a
288-step horizon, re-optimization at every interval, and execution of only
the first optimized action.
\end{frame}

% ----------------------------------------------------------------------
\begin{frame}{Water-demand calculation}
The hourly Richmond demand multipliers are converted into five-minute
values using
\[
  d_{m,\,5\mathrm{min}}=\operatorname{repelem}
  \!\left(d_{m,\,\mathrm{hourly}},12\right).
\]
The water demand during interval $k$ is
\[
  d(k)=d_{\mathrm{base}}d_m(k)=0.025d_m(k).
\]

\textbf{Where:}
\begin{itemize}
  \item $d(k)$: water demand during interval $k$, in \si{\metre\cubed\per\second};
  \item $d_{\mathrm{base}}=0.025\ \si{\metre\cubed\per\second}$:
        base demand, equivalent to \SI{25}{\litre\per\second};
  \item $d_m(k)$: dimensionless Richmond demand multiplier.
\end{itemize}

For example, if $d_m(k)=1.61$, then
$d(k)=0.025\times1.61=0.04025\ \si{\metre\cubed\per\second}
=\SI{40.25}{\litre\per\second}$.
\end{frame}

% ----------------------------------------------------------------------
\begin{frame}{Three-pump binary decisions}
For pump $j\in\{1,2,3\}$ and prediction step $k$,
\[
  u_j(k)\in\{0,1\},
  \qquad 0=\text{off},\quad 1=\text{on}.
\]
The synthetic pump parameters used by Basic 1.3 are
\[
  \bm q=\begin{bmatrix}0.050&0.030&0.020\end{bmatrix}
  \ \si{\metre\cubed\per\second},
  \qquad
  \bm P=\begin{bmatrix}43&27&19\end{bmatrix}\ \si{\kilo\watt}.
\]
Therefore, the total requested flow and power are
\[
  q(k)=\sum_{j=1}^{3}q_j u_j(k),
  \qquad
  P(k)=\sum_{j=1}^{3}P_j u_j(k).
\]

\begin{itemize}
  \item The pumps have fixed flow and fixed power; there is no continuous
        pump-speed variable.
  \item All $2^3=8$ on/off combinations are available because
        $\sum_j u_j(k)\leq3$.
  \item These values are synthetic test parameters, not measured Richmond
        pump data.
\end{itemize}
\end{frame}

% ----------------------------------------------------------------------
\begin{frame}{All eight operating modes}
\small
\begin{center}
\begin{tabular}{c l c r r r r}
\toprule
Mode & Pump combination & $[u_1\ u_2\ u_3]$ &
Flow (L/s) & Power (kW) & Energy/5 min (kWh) & PE100 \\
\midrule
0 & All off       & $[0\ 0\ 0]$ &   0 &  0 & 0.000000 & -- \\
1 & P1            & $[1\ 0\ 0]$ &  50 & 43 & 3.583333 & 86.0000 \\
2 & P2            & $[0\ 1\ 0]$ &  30 & 27 & 2.250000 & 90.0000 \\
3 & P3            & $[0\ 0\ 1]$ &  20 & 19 & 1.583333 & 95.0000 \\
4 & P1+P2         & $[1\ 1\ 0]$ &  80 & 70 & 5.833333 & 87.5000 \\
5 & P1+P3         & $[1\ 0\ 1]$ &  70 & 62 & 5.166667 & 88.5714 \\
6 & P2+P3         & $[0\ 1\ 1]$ &  50 & 46 & 3.833333 & 92.0000 \\
7 & P1+P2+P3      & $[1\ 1\ 1]$ & 100 & 89 & 7.416667 & 89.0000 \\
\bottomrule
\end{tabular}
\end{center}

\small
The five-minute electrical energy is
\[
  E_{\mathrm{step}}=P\frac{300}{3600}=\frac{P}{12}
  \quad\text{in kWh}.
\]
PE100 is an efficiency-display metric,
$\mathrm{PE100}=100P/q_{\mathrm{L/s}}$ in
$\mathrm{kW/(100\,L/s)}$.  It is \alert{not} used to calculate the
electricity bill.
\end{frame}

% ----------------------------------------------------------------------
\begin{frame}{Tank-level dynamics}
The tank level is updated using the actual net inflow:
\[
  x(k+1)=x(k)+\frac{\Delta t}{A}\left[q(k)-d(k)\right].
\]
Substituting $\Delta t=\SI{300}{\second}$, $A=\SI{1000}{\metre\squared}$,
and the three-pump flow equation gives
\[
  x(k+1)=x(k)+0.3\left[
  0.050u_1(k)+0.030u_2(k)+0.020u_3(k)-d(k)\right].
\]

\textbf{Where:}
\begin{itemize}
  \item $x(k)$: tank level at the beginning of interval $k$, in \si{\metre};
  \item $x(k+1)$: level after one five-minute interval, in \si{\metre};
  \item $A=\SI{1000}{\metre\squared}$: tank cross-sectional area;
  \item $q(k)$: combined pump inflow, in \si{\metre\cubed\per\second};
  \item $d(k)$: water demand, in \si{\metre\cubed\per\second}.
\end{itemize}

The level is a state determined by mass balance; it is not an independent
control input.
\end{frame}

% ----------------------------------------------------------------------
\begin{frame}{MILP constraints}
At every step of the prediction horizon, the controller enforces
\[
  u_j(k)\in\{0,1\},\qquad
  1.40\leq x(k)\leq3.37,\qquad
  \sum_{j=1}^{3}u_j(k)\leq3.
\]
The tank dynamics must also hold exactly:
\[
  x(k+1)-x(k)=\frac{300}{1000}\left[
  \sum_{j=1}^{3}q_j u_j(k)-d(k)\right].
\]

\begin{itemize}
  \item Initial level: $x(0)=\SI{2.30}{\metre}$.
  \item Prediction horizon: $N=288$ five-minute intervals, or 24 hours.
  \item No terminal-level equality is imposed.
  \item No flow-tracking penalty is added.
  \item Historical prices are used as perfect-forecast inputs to test the
        controller, not to validate a price predictor.
\end{itemize}
\end{frame}

% ----------------------------------------------------------------------
\begin{frame}{Per-pump switching transitions}
Each pump has a non-negative switching variable $s_j(k)$.  For the first
predicted interval,
\[
  s_j(1)\geq u_j(1)-u_{j,\mathrm{previous}},\qquad
  s_j(1)\geq u_{j,\mathrm{previous}}-u_j(1).
\]
For later intervals,
\[
  s_j(k)\geq u_j(k)-u_j(k-1),\qquad
  s_j(k)\geq u_j(k-1)-u_j(k).
\]
Because $s_j(k)$ is minimized with a positive coefficient,
\[
  s_j(k)=\left|u_j(k)-u_j(k-1)\right|
\]
at the optimum.

\begin{itemize}
  \item Switching cost: $c_s=\SI{0.20}{\AUD}$ per transition per pump.
  \item An on transition and an off transition are counted separately.
  \item Changing from P1+P2 to P1+P3 produces two transitions:
        P2 turns off and P3 turns on.
  \item The previous state comes from the pumps actually executed in the
        preceding closed-loop interval.
\end{itemize}
\end{frame}

% ----------------------------------------------------------------------
\begin{frame}{Electricity-cost conversion}
The electricity price $p(k)$ is in \si{\AUD\per\mega\watt\hour}, whereas
pump power $P(k)$ is in \si{\kilo\watt}.  The five-minute energy cost is
\[
  C_{\mathrm{energy}}(k)
  =p(k)P(k)\frac{300}{3600}\frac{1}{1000}.
\]

\textbf{Unit check:}
\[
  \frac{\mathrm{AUD}}{\mathrm{MWh}}
  \times\mathrm{kW}
  \times\mathrm{h}
  \times\frac{\mathrm{MW}}{1000\,\mathrm{kW}}
  =\mathrm{AUD}.
\]

For example, one five-minute P1 interval uses
\[
  E_{1,\mathrm{step}}=43\times\frac{300}{3600}
  =\SI{3.583333}{\kilo\watt\hour}
  =\SI{0.003583333}{\mega\watt\hour}.
\]

The factor $1/1000$ must be retained.  PE100 is used only for efficiency
display; the electricity bill uses the actual powers 43, 27, and
\SI{19}{\kilo\watt}.
\end{frame}

% ----------------------------------------------------------------------
\begin{frame}{Three-pump MPC objective function}
At every sampling instant, the MILP minimizes predicted electricity and
per-pump switching costs over $N=288$ future intervals:
\[
  \min_{u,s}\ J=
  \sum_{k=1}^{N}\left[
    p(k)\left(\sum_{j=1}^{3}P_j u_j(k)\right)
    \frac{300}{3600}\frac{1}{1000}
    +c_s\sum_{j=1}^{3}s_j(k)
  \right].
\]

\textbf{Where:}
\begin{itemize}
  \item $J$: predicted operating cost over the 24-hour horizon, in AUD;
  \item $p(k)$: predicted electricity price, in \si{\AUD\per\mega\watt\hour};
  \item $P_j$: fixed power of pump $j$, in \si{\kilo\watt};
  \item $u_j(k)$: planned binary state of pump $j$;
  \item $s_j(k)$: switching transition of pump $j$;
  \item $c_s=\SI{0.20}{\AUD}$ per transition per pump.
\end{itemize}

There is no separate penalty on pumped volume or final tank level.
\end{frame}

% ----------------------------------------------------------------------
\begin{frame}{Rolling closed-loop MPC operation}
At closed-loop interval $t$, Basic 1.3 performs the following sequence:
\begin{enumerate}
  \item read the current measured tank level $x(t)$ and the previously
        executed pump state $\bm u_{\mathrm{actual}}(t-1)$;
  \item build a 288-step price and demand prediction window;
  \item solve the three-pump MILP using \code{intlinprog};
  \item obtain the planned matrix
        $\bm U^\star\in\{0,1\}^{288\times3}$;
  \item execute only the first row,
        $\bm u_{\mathrm{request}}(t)=\bm U^\star(1,:)$;
  \item obtain actual pump states, total flow, and total power from the
        execution interface;
  \item update the tank with actual flow and calculate cost using actual
        power;
  \item advance five minutes and solve a new 288-step MILP.
\end{enumerate}

Thus a 24-hour closed-loop test solves 288 separate MILPs; it is not one
open-loop 24-hour schedule.
\end{frame}

% ----------------------------------------------------------------------
\begin{frame}{Absolute cost versus cost per cubic metre}
The historical comparison produced the following realised results:

\begin{center}
\begin{tabular}{lrr}
\toprule
Metric & Multi-pump & P1-only \\
\midrule
Pumped volume & \SI{1830}{\metre\cubed} & \SI{1590}{\metre\cubed} \\
Final tank level & \SI{1.9781}{\metre} & \SI{1.7381}{\metre} \\
Absolute total cost & \SI{29.7296}{\AUD} & \SI{27.7452}{\AUD} \\
Cost per pumped volume & \SI{0.01625}{\AUD\per\metre\cubed}
                       & \SI{0.01745}{\AUD\per\metre\cubed} \\
\bottomrule
\end{tabular}
\end{center}

\begin{itemize}
  \item \textbf{Absolute cost:} P1-only is cheaper because it pumps
        \SI{240}{\metre\cubed} less water during the displayed day.
  \item \textbf{Cost per cubic metre:} P1-only is more expensive per unit
        of pumped water.  Multi-pump is approximately $6.9\%$ cheaper on
        this metric.
  \item The final inventories are different, so absolute cost alone is not
        an equal-service economic comparison.
\end{itemize}
\end{frame}

% ----------------------------------------------------------------------
\begin{frame}{Why the two conclusions are both possible}
The volume difference is exactly explained by final tank inventory:
\[
  \Delta V=A\Delta x
  =1000(1.9781-1.7381)=\SI{240}{\metre\cubed}.
\]
This equals the pumping difference $1830-1590=\SI{240}{\metre\cubed}$:
the extra water remained in the tank and was not extra demand.

\begin{itemize}
  \item P1-only has the lower \emph{total daily bill}.
  \item Multi-pump has the lower \emph{cost per pumped cubic metre}, because
        more pumping is shifted into cheaper price intervals.
  \item The lower relative quantity is \textbf{AUD/m$^3$}, not
        \textbf{kWh/m$^3$}; therefore ``relative consumption'' is ambiguous.
\end{itemize}

The electrical intensities are
\[
  \text{Multi}=0.2448\ \mathrm{kWh/m^3},\qquad
  \text{P1-only}=0.2389\ \mathrm{kWh/m^3}.
\]
\textbf{Conclusion:} the economic advantage per cubic metre comes from
electricity-price timing, not from lower physical energy use.
\end{frame}

% ----------------------------------------------------------------------
\begin{frame}{Interpretation and fair comparison}
The multi-pump feasible set contains the P1-only actions, but this does not
guarantee a lower truncated 24-hour realised bill when the two trajectories
finish with different stored water.

\begin{itemize}
  \item Every rolling optimization looks 24 hours ahead, including times
        beyond the end of the displayed day.
  \item With no terminal-level equality, the controller may use a cheap
        period today to store water for future demand.
  \item The reported comparison is fair in price, demand, initial level,
        constraints, horizon, and simulation duration, but not in final
        inventory.
\end{itemize}

For a strict economic comparison, report all three quantities together:
\[
  \boxed{\text{total cost}},\qquad
  \boxed{\text{cost per pumped }\mathrm{m^3}},\qquad
  \boxed{\text{final stored volume}}.
\]
Alternatively, run a separate analysis in which both controllers are
required to finish at the same tank level.  This is an evaluation choice;
it is not part of the current Basic 1.3 controller specification.
\end{frame}

\end{document}
